\section*{Chapter \ref{ch:ll:protocols-i-fix} --- Protocols I: FIX}

\begin{solution}{exo:ll:protocols-i-fix:1}
The body is \texttt{35=0|} (5 bytes), \texttt{49=A|} (5), \texttt{56=B|} (5), \texttt{34=2|} (5) and the sending time field (25):
45 bytes, so \texttt{9=45}. The bytes before \texttt{10=} sum to a number whose remainder modulo 256 is 75: \texttt{10=075}, as
\texttt{firm\_fixengine.encode} confirms for sequence number 2 at 14:30:00.
\end{solution}

\begin{solution}{exo:ll:protocols-i-fix:2}
A test request after 1.2 intervals of silence, at 10:00:36; if no \omterm{def:ll:protocols-i-fix:session}{heartbeat} carrying its identifier arrives within one more
interval, the session declares the connection dead at 10:01:06, 66 seconds after the last message.
\end{solution}

\begin{solution}{exo:ll:protocols-i-fix:3}
A ResendRequest with BeginSeqNo 52 and EndSeqNo 0; it holds message 57. The broker resends 52--56 (and everything after, as
EndSeqNo 0 asks): application messages with PossDupFlag and OrigSendingTime, administrative ones replaced by
SequenceReset-GapFill; then 57 is delivered from the queue, and the resent copy of 57 is ignored as a duplicate.
\end{solution}

\begin{solution}{exo:ll:protocols-i-fix:4}
Swapping the sender and target identifiers of the \omterm{def:ll:protocols-i-fix:session}{heartbeat} of exercise~1 gives \texttt{49=B|56=A|} instead of
\texttt{49=A|56=B|}: the same bytes in another order, so the same length and the same checksum, 075. The checksum only has to
catch framing errors and single-byte damage; TCP's own checksums protect the bytes in transit, and the session layer rejects
a message addressed to the wrong party by its identifiers.
\end{solution}

\begin{solution}{exo:ll:protocols-i-fix:5}
They describe the link at the moment they were sent: a \omterm{def:ll:protocols-i-fix:session}{heartbeat} says ``I am alive now'', a test request asks for an answer
now, a \omterm{def:ll:protocols-i-fix:session}{resend request} asks for messages that may since have been sent. Resent later they would be meaningless or harmful: a
stale test request would demand an answer, a stale \omterm{def:ll:protocols-i-fix:session}{resend request} would trigger a second resend. \omterm{def:ll:protocols-i-fix:session}{Gap fill} moves the sequence
past them without replaying them.
\end{solution}

\begin{solution}{exo:ll:protocols-i-fix:6}
On the committed measurement: $200\,000 \times \qty{84}{\nano\second} \approx 1.7\%$ of a core with the zero-copy view,
$200\,000 \times \qty{892}{\nano\second} \approx 18\%$ with the map-based parser, and $200\,000 \times
\qty{5.3}{\micro\second} \approx 107\%$ with the Python reference, which cannot keep up at all.
\end{solution}

\begin{solution}{exo:ll:protocols-i-fix:7}
In \texttt{on\_bytes}, catch the decoding error of kind \texttt{tag} before the sequence check, read MsgSeqNum from the raw bytes
if it is present, and emit \texttt{35=3} with \texttt{45=<that number>} and a text; then advance the expected number (the
message was received, only not processed). Do the same in the C++ \texttt{on\_bytes} on \texttt{Error::tag}, add a malformed
message to \texttt{make\_fix\_fixtures.py}, regenerate the expected trace and run both tests.
\end{solution}

\begin{solution}{exo:ll:protocols-i-fix:8}
Messages sent while the connection was down (or in flight when it broke) are never delivered: an execution report for a fill
is lost, and the firm's position is wrong. Resetting is acceptable only with another way to recover the state of every order,
such as a drop copy or order-status queries after each reconnection, as venues that support no \omterm{def:ll:protocols-i-fix:session}{resend requests} expect.
\end{solution}

\begin{solution}{pb:ll:protocols-i-fix:1}
\begin{enumerate}
\item About 66 seconds: a test request after 36 seconds of silence and one more 30-second interval for its answer.
\item 13\,227 execution reports on the fixed seed, and no \omterm{def:ll:protocols-i-fix:session}{heartbeat}: at 200 messages a second the broker is never idle for 30
seconds.
\item $13\,227 \times 205 = 2\,711\,535$ bytes, about \qty{21.7}{\milli\second} at 1~Gb/s.
\item Two round trips (logon, \omterm{def:ll:protocols-i-fix:session}{resend request}), the bytes, and the parsing: about \qty{24.8}{\milli\second} with the C++
engine and about \qty{94}{\milli\second} with the Python reference.
\item 5\,776 messages, of which 38 are application messages and the rest \omterm{def:ll:protocols-i-fix:session}{heartbeats}.
\item With \omterm{def:ll:protocols-i-fix:session}{gap fill} 77: the 38 application messages and 39 \omterm{def:ll:protocols-i-fix:session}{gap fills}, one per run of \omterm{def:ll:protocols-i-fix:session}{heartbeats}. Without, 5\,776.
\item About \qty{2.5}{\milli\second} with \omterm{def:ll:protocols-i-fix:session}{gap fill} against \qty{42}{\milli\second} without.
\item There are no administrative messages to replace.
\item They refer to the link at the time they were sent (exercise~5).
\item Check whether it has already been processed, by its execution identifier, and book it only once.
\item From the venue's drop copy or by querying the status of every open order after reconnecting, and reconcile with the
firm's own records before trading again.
\item Its next outgoing and expected incoming sequence numbers, and the messages it sent (to answer \omterm{def:ll:protocols-i-fix:session}{resend requests}), written
before or as they are sent.
\item \textbf{Named result.} A busy 66-second outage at 200 messages a second leaves 13\,227 messages to resend, about
\qty{25}{\milli\second} to resynchronise with the C++ engine and \qty{94}{\milli\second} with the Python reference, \omterm{def:ll:protocols-i-fix:session}{gap fill} or
not; a quiet 48-hour weekend leaves 5\,776 messages, which \omterm{def:ll:protocols-i-fix:session}{gap fill} reduces to 77, from \qty{42}{\milli\second} to
\qty{2.5}{\milli\second}.
\item Yes: detection in 2.2 seconds leaves about 450 messages instead of 13\,227; the cost is a \omterm{def:ll:protocols-i-fix:session}{heartbeat} a second each way
when idle, and false alarms whenever a peer is briefly slow.
\item Messages older than the store cannot be resent: the broker must gap-fill or reset over them, and the lost execution
reports must be recovered by reconciliation.
\item It should update positions and order states from them at once, but judge them as past events: prices and the book have
moved since their original sending time (122).
\item With golden conversations (lost messages, duplicates, \omterm{def:ll:protocols-i-fix:session}{resend requests} in both directions, resets), as the build does,
and against the exchange simulator of One Quant Book~10.
\item Yes: 200 messages a second at \qty{5.3}{\micro\second} is under 0.1\% of a core; it is the resend burst that it takes
tens of milliseconds to absorb.
\item The busy case: its resend is the large one, and it happens when markets are moving.
\item To keep two firms' views of their conversation identical, message for message, whatever the network does.
\end{enumerate}
\end{solution}

\begin{solution}{iq:ll:protocols-i-fix:1}
Fields \texttt{tag=value} separated by SOH: BeginString, BodyLength and MsgType first, the rest of the header, the body, and
CheckSum last. The receiver reads the first two fields; BodyLength gives the length up to the checksum field, which is always
seven bytes (\texttt{10=nnn} and SOH).
\iqlookfor{body length for framing, checksum as the terminator.}
\end{solution}

\begin{solution}{iq:ll:protocols-i-fix:2}
The receiver sees a sequence number above the expected one, sends a ResendRequest from the expected number (EndSeqNo 0), and
holds the later messages. The sender resends application messages with PossDupFlag and OrigSendingTime, and gap-fills
administrative ones with SequenceReset. The receiver processes them in order, releases the held messages, and ignores
duplicates.
\iqlookfor{hold, resend, \omterm{def:ll:protocols-i-fix:session}{gap fill}, duplicates.}
\end{solution}

\begin{solution}{iq:ll:protocols-i-fix:3}
Find every delimiter (vectorised), parse each tag in place, and record (tag, offset, length) in a fixed array; values are
views into the receive buffer, valid as long as it is. Validate the framing fields, body length and checksum in the same pass.
Return the view, or an error kind.
\iqlookfor{views, a fixed array, and validation in one pass.}
\end{solution}

\begin{solution}{iq:ll:protocols-i-fix:4}
To tell silence from a dead connection: a side idle for the interval sends a \omterm{def:ll:protocols-i-fix:session}{heartbeat}; a side that hears nothing sends a test
request and gives up if it gets no answer. A shorter interval detects failures sooner, so fewer messages go into the void, at
the price of more traffic and of false alarms when a peer pauses.
\iqlookfor{detection time against traffic and false alarms.}
\end{solution}

\begin{solution}{iq:ll:protocols-i-fix:5}
Fixed-offset binary fields need no search, no number parsing and less bandwidth, and fixed sizes simplify framing: tens of
nanoseconds per message, predictable. They give up readability, the common dictionary shared across venues, and flexibility:
each venue has its own layouts and versions.
\iqlookfor{parse cost and determinism against interoperability.}
\end{solution}

\begin{solution}{iq:ll:protocols-i-fix:6}
The engine lost its outgoing sequence number (not persisted, or restored from an old copy) and restarted below the broker's
expected number. Restore the persisted numbers and store; if they are lost, agree a sequence reset with the broker out of band,
then reconcile every order's state from the broker's records or a drop copy before trading, booking fills by execution
identifier so that none is counted twice.
\iqlookfor{persisted state, a controlled reset, and reconciliation.}
\end{solution}
